\section{Kết luận}

\subsection{Tóm tắt kết quả đạt được}

Trong khuôn khổ đồ án, nhóm đã hoàn thành quy trình Data Fitting theo đúng yêu cầu:
từ nền tảng lý thuyết OLS, cài đặt các thuật toán hồi quy từ công thức, đến xây dựng
pipeline tiền xử lý và so sánh mô hình trên dữ liệu ngoại hành tinh thực tế. Các kết quả
trọng tâm đạt được gồm:

\begin{itemize}[itemsep=0.15cm, topsep=0.1cm]
    \item \textbf{Hoàn thiện nền tảng Data Fitting và OLS:} Nhóm đã trình bày bài toán hồi quy
          tuyến tính, hàm mất mát RSS, nghiệm OLS, Hat Matrix, định lý Gauss--Markov, kiểm định
          hệ số, $R^2$, $R^2$ hiệu chỉnh và phân tích phần dư. Các nội dung lý thuyết được minh
          họa bằng dữ liệu giả lập và kiểm chứng với NumPy/sklearn để đảm bảo phần cài đặt từ
          đầu cho kết quả nhất quán.

    \item \textbf{Cài đặt đầy đủ các công cụ mô hình hóa cần thiết:} Các hàm OLS, Ridge, Lasso,
          VIF, k-fold cross-validation, kiểm định hệ số, residual diagnostics và model metrics
          đã được hiện thực bằng Python. Phần cài đặt này không chỉ phục vụ minh họa lý thuyết
          ở Phần 1 mà còn được tái sử dụng trực tiếp trong Phần 2 để huấn luyện và đánh giá mô
          hình trên dữ liệu thật.

    \item \textbf{Chuẩn hóa bộ dữ liệu ngoại hành tinh cho bài toán hồi quy:} Dữ liệu chuẩn
          được đưa vào từ \texttt{part2/data/data.csv} với 4636 dòng và 320 cột gốc của NASA
          Exoplanet Archive. Từ tập này, nhóm áp dụng schema filter để giữ 16 biến vật lý phù
          hợp với bài toán dự đoán bán kính hành tinh, sau đó giới hạn miền nghiên cứu theo
          Fulton Gap: \texttt{pl\_rade}$\le1.6$ và \texttt{pl\_bmasse}$\le10$. Tập mô hình hóa
          cuối cùng còn 1084 mẫu thuộc nhóm hành tinh đá/Super-Earth.

    \item \textbf{Xây dựng pipeline tiền xử lý có kiểm soát leakage:} Pipeline gồm
          log-transform, Winsorization, MICE imputation, VIF filtering và standardization. Các
          tham số tiền xử lý chỉ được fit trên train rồi mới áp dụng cho test, giúp tránh rò rỉ
          dữ liệu đánh giá. Việc dùng MICE thay cho xóa dòng hoặc điền trung bình giúp giữ lại
          cấu trúc tương quan giữa các nhóm biến thiên văn trong bối cảnh dữ liệu thiếu không
          đồng đều.

    \item \textbf{Thanh lọc đặc trưng bằng kết hợp kiến thức vật lý và thống kê:} Trước pipeline,
          nhóm loại các cột metadata, sai số đo lường, biến gây rò rỉ khái niệm, đại lượng dẫn
          xuất và các biến trùng thông tin. Trong pipeline, VIF tiếp tục loại
          \texttt{log\_pl\_orbper} và \texttt{log\_pl\_eqt}, đưa ma trận thiết kế cuối cùng về
          11 feature. Cách chọn biến này giúp mô hình tránh đa cộng tuyến nhưng vẫn giữ các tín
          hiệu vật lý chính của hệ hành tinh--sao chủ.

    \item \textbf{So sánh có hệ thống các mô hình hồi quy:} Nhóm đã huấn luyện và đánh giá
          OLS full, OLS selected, Ridge và Lasso bằng train/test split cùng 5-fold
          cross-validation để chọn siêu tham số. Trong lần chạy chính thức, Lasso là mô hình
          tuyến tính tốt nhất trên test set với $R^2=0.748403$, RMSE $=0.056872$ và
          MAE $=0.037895$ trên thang log target. Khi biến đổi ngược về đơn vị bán kính Trái Đất,
          Lasso đạt $R^2=0.712850$, RMSE $=0.130278$, MAE $=0.085171$ và MAPE $=7.032\%$.

    \item \textbf{Giải thích kết quả mô hình thay vì chỉ báo cáo số liệu:} Sự khác biệt giữa
          Lasso, OLS full và Ridge rất nhỏ vì log-transform, Winsorization và VIF filtering đã
          xử lý phần lớn lệch phân phối và đa cộng tuyến trước khi mô hình học. Với
          $\lambda^*_{\mathrm{Lasso}}=0.1$, Lasso không triệt tiêu feature nào ($0/11$ hệ số
          bằng 0), cho thấy regularization chủ yếu đóng vai trò co nhẹ hệ số hơn là chọn biến
          thưa trong bộ dữ liệu sau tiền xử lý.

    \item \textbf{Hoàn thành đánh giá chẩn đoán và thảo luận giới hạn:} Báo cáo đã bổ sung
          residual plots, Shapiro-Wilk, Breusch-Pagan, Cook's Distance, phân tích các điểm
          influential và thảo luận selection bias của \texttt{log\_sy\_dist}. Nhóm không xóa
          các điểm influential một cách cơ học vì đây là dữ liệu quan trắc thiên văn thật; thay
          vào đó, nhóm ghi nhận chúng như một giới hạn thực nghiệm của mô hình tuyến tính.

    \item \textbf{Thử nghiệm kỹ thuật nâng cao bằng Kernel Ridge Regression:} Kernel Ridge
          Regression với RBF kernel được triển khai đúng công thức kernel trick trong đề bài.
          Khi chạy trên toàn bộ 867 mẫu train sau tiền xử lý, mô hình đạt $R^2=0.855688$,
          RMSE $=0.043072$ và MAE $=0.029685$ trên test set, cao hơn nhóm mô hình tuyến tính.
          So với Lasso tốt nhất trong nhóm tuyến tính, Kernel Ridge tăng Test $R^2$ thêm
          0.107285 điểm, giảm RMSE từ 0.056872 xuống 0.043072 và giảm phần phương sai chưa
          giải thích được khoảng 42.6\%. Kết quả này gợi ý dữ liệu có thành phần phi tuyến
          đáng kể, nhưng mô hình kernel vẫn được đặt ở vai trò bonus vì khó diễn giải hệ số
          và chi phí ma trận Gram/giải hệ tăng nhanh theo kích thước mẫu.
\end{itemize}

\subsection{Bài học rút ra}

Quá trình thực hiện đồ án giúp nhóm rút ra một số bài học quan trọng về xây dựng mô hình hồi quy
trên dữ liệu thực:

\begin{itemize}[itemsep=0.15cm]
    \item \textbf{Mô hình tốt bắt đầu từ dữ liệu đúng:} Với dữ liệu ngoại hành tinh, việc chọn
          đúng miền vật lý quan trọng không kém thuật toán. Nếu trộn hành tinh khí khổng lồ với
          hành tinh đá, mô hình tuyến tính phải học nhiều cơ chế vật lý khác nhau và dễ sinh sai
          lệch phương sai.

    \item \textbf{Tiền xử lý phải được fit đúng quy trình:} Log-transform, Winsorization, MICE,
          VIF và standardization đều cần tách rõ train/test. Chỉ một bước fit nhầm trên toàn bộ
          dữ liệu cũng có thể làm test set mất ý nghĩa đánh giá độc lập.

    \item \textbf{Regularization không phải lúc nào cũng tạo cải thiện lớn:} Ridge và Lasso phát
          huy mạnh khi dữ liệu còn đa cộng tuyến hoặc nhiều biến nhiễu. Sau khi đã lọc VIF và
          loại các biến rò rỉ/dẫn xuất, lợi thế của regularization giảm đi, nên kết quả gần OLS
          là hợp lý.

    \item \textbf{Chỉ số đánh giá cần đi kèm diễn giải ngữ cảnh:} $R^2_{\mathrm{test}}\approx
          0.748$ cho thấy mô hình giải thích được phần lớn phương sai trên thang log, nhưng vẫn
          còn khoảng 25\% phương sai chưa giải thích được. Phần còn lại có thể đến từ quan hệ
          phi tuyến, sai số quan trắc, selection bias và các biến vật lý chưa có trong dữ liệu.

    \item \textbf{Không xử lý ngoại lai một cách máy móc:} Cook's Distance giúp phát hiện điểm có
          ảnh hưởng lớn, nhưng trong dữ liệu khoa học thật, điểm lạ không nhất thiết là lỗi. Vì
          vậy nhóm ưu tiên giảm ảnh hưởng bằng log-transform/Winsorization và thảo luận minh bạch
          thay vì xóa dữ liệu tùy tiện.
\end{itemize}

\subsection{Khó khăn và hướng mở rộng}

Khó khăn lớn nhất của đồ án nằm ở việc cân bằng giữa yêu cầu toán học của mô hình tuyến tính và
bản chất phức tạp của dữ liệu thiên văn. Dữ liệu có nhiều missing values, nhiều cột là metadata
hoặc đại lượng dẫn xuất, đồng thời các biến vật lý thường liên hệ với nhau theo quy luật hàm mũ,
hàm lũy thừa hoặc cơ chế quan trắc không đồng đều. Vì vậy, nhóm phải kết hợp kiến thức miền với
các công cụ thống kê như VIF, residual analysis và cross-validation để đưa ra quyết định có căn cứ.

Trong các hướng mở rộng tiếp theo, nhóm có thể thử các mô hình phi tuyến trên toàn bộ tập train
với chiến lược tối ưu bộ nhớ tốt hơn, chẳng hạn Kernel Ridge theo mini-batch/approximate kernel,
Gaussian Process xấp xỉ hoặc các mô hình cây tăng cường. Ngoài ra, có thể tách riêng các nhóm
hành tinh theo miền vật lý nhỏ hơn, bổ sung uncertainty-aware modeling cho sai số đo lường, và
kiểm tra độ ổn định của kết quả qua nhiều seed train/test split.
